\clearpage\thispagestyle{empty}
\pdfbookmark[0]{Sources for Gauss to von Lindenau, 30 November 1808}{lindenau1808nov-sources}
\begin{center}\textbf{Sources for Gauss to von Lindenau, 30 November 1808}\end{center}
\noindent This AI translation follows Carl Friedrich Gauss’s German letter in \textit{Werke}, vol.~VI (Göttingen, 1874), pp.~141--145, including its handwritten remark. The \href{https://archive.org/details/bub_gb_caMAAAAAMAAJ/page/n91/mode/1up}{first printing}, article V of the January 1809 issue of Franz Xaver von Zach’s \textit{Monatliche Correspondenz zur Beförderung der Erd- und Himmels-Kunde}, vol.~19, pp.~85--92, contains a note signed v.~L. on p.~86 and Gauss’s list of misprints in his October paper on p.~90. \href{https://gdz.sub.uni-goettingen.de/id/PPN236007467}{Collected German text}.
\par\smallskip\noindent Joseph-Jérôme de Lalande describes Cagnoli’s method and its numerical example in \href{https://doi.org/10.3931/e-rara-9241}{\textit{Astronomie}}, 3rd ed., vol.~III (Paris, 1792), arts.~3259--3261. The logarithmic tables Gauss mentions are Michael Taylor’s \href{https://doi.org/10.3931/e-rara-9132}{\textit{Tables of Logarithms}} (London, 1792), with preface and instructions by Nevil Maskelyne. Franz Brünnow treats Gauss’s method with his own example in \href{https://archive.org/details/bub_gb_sAA5AAAAMAAJ/page/n306/mode/1up}{\textit{Lehrbuch der sphärischen Astronomie}} (Berlin, 1851), §20, and in his English translation from the second German edition, \href{https://archive.org/details/sphericalastro00brhurich/page/n328/mode/1up}{\textit{Spherical Astronomy}} (London, 1865), §20.
\par\smallskip\noindent Translation and typesetting produced by AI.
\clearpage
